% Corrected French for printed p.143.
% Correction: restore the missing quotient slash in F(M/A)/Im(F(M)).
\SourcePage{143}

\LayerTarget{tohoku:remark:satellites-small-skeleton:p2}
\noindent abéliens définis sur $k$ dont la composante connexe de l'élément
neutre est une ``variété abélienne''. (On est obligé de supposer la
caractéristique nulle, sinon un homomorphisme bijectif ne serait pas
nécessairement un isomorphisme). Indiquons seulement qu'on démontre le
résultat énoncé ci-dessus en construisant $S^1F(A)$ comme limite inductive des
objets $F(M/A)/\mathrm{Im}(F(M))$ pour ``tous'' les monomorphismes
$A\to M$ dans $\mathbf{C}$, préordonnés en disant que
$A\to M$ est ``majoré'' par $A\to M'$ si on peut
trouver un morphisme $M\to M'$ induisant l'identité sur $A$.

\LayerTarget{tohoku:section:2.3}
\indent\textbf{2.3. Foncteurs dérivés.} Soient $\mathbf{C}$ et $\mathbf{C}'$ deux catégories
abéliennes. La théorie des foncteurs dérivés d'un foncteur additif $F$ de $\mathbf{C}$
dans $\mathbf{C}'$ se développe comme dans [6, Chap. V], à cela près qu'il faut supposer
suivant les cas que tout objet $A\in \mathbf{C}$ est isomorphe à un sous-truc d'un objet
injectif, resp. à un truc quotient d'un objet projectif, ou les deux.
\LayerTarget{tohoku:construction:right-derived-functors:p1}
Ainsi, pour pouvoir définir les foncteurs dérivés \textit{droits} d'un foncteur
\textit{covariant} ou les foncteurs dérivés \textit{gauches} d'un foncteur
\textit{contravariant}, il faut supposer que tout objet $A\in \mathbf{C}$ est isomorphe
à un sous-truc d'un objet injectif, d'où on conclut en effet que tout $A\in \mathbf{C}$
admet une résolution injective\textsuperscript{3)}\SourceFootnote{3}{Une \textit{résolution droite} d'un
$A\in \mathbf{C}$ est par définition un complexe $C$ à degrés positifs, muni d'un
``homomorphisme d'augmentation'' $A\to C$ ($A$ étant considéré
comme un complexe réduit au degré 0), de façon que la suite
$0\to A\to C^0\to C^1\to
\ldots\ldots$ soit exacte. On appelle \textit{résolution injective} de $A$ une
résolution $C$ de $A$ telle que les $C^i$ soient des objets injectifs. Les
\textit{résolutions gauches} de $A$, et en particulier les \textit{résolutions
projectives} de $A$, se définissent dualement.} :
$0\to A\to C^0\to C^1\to
\ldots\,.$, d'où la définition des
$R^iF(A)=H^i(F(C))$ (où $C$ désigne le complexe des $C^i$). Si on veut définir
des foncteurs dérivés gauches d'un foncteur covariant ou des foncteurs dérivés
droits d'un foncteur contravariant, il faut de même supposer que tout objet
$A\in \mathbf{C}$ est isomorphe à un quotient d'un objet projectif. Et enfin, pour
définir des foncteurs dérivés d'un foncteur mixte en plusieurs variables, il
faut faire sur la catégorie de chaque variable l'hypothèse appropriée. A cela
près, l'exposé de [6] se transporte tel quel.
\LayerTarget{tohoku:morphism:satellite-to-derived}
En particulier, si $F$ est par exemple covariant et si on peut former les
foncteurs dérivés droits $R^iF$, alors (posant $R^iF=0$ pour $i<0$)
$RF=(R^iF)$ est un foncteur cohomologique (dit \textit{foncteur cohomologique
dérivé droit} de $F$), et on a un morphisme canonique de $\partial$-foncteurs
à degrés positifs $SF\to RF$ (où $SF=(S^iF)$ est le
$\partial$-foncteur universel à degrés positifs satellite de $F$, qui existe
en vertu de th. 2.2.2.) ; ce dernier est un isomorphisme si et seulement si
$F$ est exact à gauche. Notons qu'il semble que la considération des $R^iF$
n'ait guère d'intérêt que dans le cas où $F$ est exact à gauche, i. e. quand ils
coïncident avec les foncteurs satellites ; cependant la définition simultanée
des $R^iF$ par résolutions injectives est plus maniable que la définition
récurrente des $S^iF$, et en particulier se prête mieux à la construction des
suites spectrales les plus importantes (voir 2.4).

\LayerTarget{tohoku:construction:derived-bifunctor:p1}
Soient $\mathbf{C}_1,\mathbf{C}_2,\mathbf{C}'$ trois catégories abéliennes, soit $T(A,B)$ un bifoncteur
additif de $\mathbf{C}_1\times \mathbf{C}_2$ dans $\mathbf{C}'$, que nous supposerons pour fixer les idées
contravariant en $A$ et covariant en $B$. Supposons que tout objet de $\mathbf{C}_2$
soit isomorphe à un sous-truc d'un objet injectif, on peut alors construire les

\TohokuContinuation{144}
